\documentclass[a4paper,11pt]{article}
\usepackage[a4paper,margin=2.4cm]{geometry}
\usepackage{amsmath,amssymb,amsthm}
\usepackage{fontspec}
\usepackage{unicode-math}
\usepackage{graphicx}
\usepackage{booktabs}
\usepackage{hyperref}
\usepackage{xcolor}
\definecolor{neonmagenta}{HTML}{FF10F0}
\usepackage{enumitem}
\usepackage{seqsplit}
\usepackage{fvextra}
% Allow long filenames to break across lines
\newcommand{\filename}[1]{\texttt{\seqsplit{#1}}}
% fvextra Verbatim for code listings (Unicode-safe)
\newcounter{codelisting}
\newenvironment{codelisting}[2]{%
  \refstepcounter{codelisting}\label{#2}%
  \begin{center}\textit{\small Listing \thecodelisting: #1}\end{center}%
  \Verbatim[breaklines=true,breakanywhere=true,fontsize=\footnotesize,frame=single,framesep=4pt,rulecolor=\color{gray!50},fillcolor=\color{gray!5}]%
}{%
  \endVerbatim%
}
% Fonts: Fira Sans (text), Fira Mono (code), Fira Math (math) are primary.
% Where a glyph is missing from the Fira fonts, it is substituted from
% Noto Sans (text) or Noto Sans Math (math and listing symbols):
%  - U+2208 (in) is missing from Fira Mono; in text/listings it is taken
%    from Noto Sans Math. Fira Math has it, so math-mode \in is Fira.
%  - U+2090 (subscript a) is missing from Fira Mono; in text/listings it
%    is taken from Noto Sans. In math it stays a roman subscript.
%  - Math glyphs missing from Fira Math are substituted by range from
%    Noto Sans Math (range assignments added below the \setmathfont line
%    for exactly the codepoints the Fira Math coverage lacks).
\setmainfont{Fira Sans}[BoldFont={Fira Sans Bold},ItalicFont={Fira Sans Italic},BoldItalicFont={Fira Sans Bold Italic}]
\setmonofont{Fira Mono}[BoldFont={Fira Mono Bold}]
\setmathfont{Fira Math}[Scale=1.15]
% Math glyphs missing from Fira Math are substituted by range from Noto Sans Math:
% U+2216 (setminus) is absent from Fira Math.
\setmathfont{Noto Sans Math}[range={"2216},Scale=1.15]
\newfontfamily\notosans{Noto Sans}
\newfontfamily\notosansmath{Noto Sans Math}
\catcode`\∈=\active
\protected\def∈{\ifmmode\in\else{\notosansmath\char"2208}\fi}
\catcode`\ₐ=\active
\protected\defₐ{\ifmmode{}_{\mathrm{a}}\else{\notosans\char"2090}\fi}
\hypersetup{colorlinks=true,linkcolor=blue!60!black,citecolor=blue!60!black,urlcolor=blue!60!black}
\setlength{\parskip}{0.6em}
\setlength{\parindent}{0pt}
\setlength{\emergencystretch}{2em}

\theoremstyle{plain}
\newtheorem{theorem}{Theorem}[section]
\newtheorem{proposition}[theorem]{Proposition}
\newtheorem{lemma}[theorem]{Lemma}
\newtheorem{corollary}[theorem]{Corollary}

\theoremstyle{definition}
\newtheorem{definition}[theorem]{Definition}
\newtheorem{example}[theorem]{Example}

\theoremstyle{remark}
\newtheorem{remark}[theorem]{Remark}

\begin{document}

\title{\textbf{\Large Dicriticals at Infinity of Keller Maps:}\\[0.6em]
\large Theorem R and Local Obstructions\\[0.4em]
\normalsize Counting Dicriticals, the Line at Infinity, First Blow-Ups, and the $m=1$ Bridge under Its Premises}
\author{\textbf{Brandon D. McCrary}\\[0.25em]
\textit{Independent Researcher}\\[0.25em]
\href{https://orcid.org/0009-0008-2804-7165}{ORCID: 0009-0008-2804-7165}\\[0.6em]
{\small \parbox{0.92\textwidth}{\centering \textbf{Computation and LaTeX Source:}\\
{\color{neonmagenta}\hypersetup{urlcolor=neonmagenta}\href{https://github.com/brandonmccraryresearch-cloud/2D-Jacobian-Conjecture-Workspace}{\nolinkurl{https://github.com/brandonmccraryresearch-cloud/2D-Jacobian-Conjecture-Workspace}}}}}}
\date{November 2, 2026}
\maketitle

\begin{center}
\small\textbf{Revision note.} This paper supersedes the October 1, 2026 draft
\textit{No Keller Map Admits Two Dicritical Divisors: Closure of the Two-Dicritical Line
for the Planar Jacobian Conjecture}. That draft's main theorem is \textbf{false as worded},
and its claims of closure and of unconditional status are \textbf{withdrawn} here. The
corrections follow the restatement note of October 13, 2026
(\filename{branch\_two\_dicritical/RESTATEMENT.md}, unpublished working note): the corrected
target is Theorem~R, stated as \textbf{open} in \S\ref{sec:statement}, and every result below
is labelled with its status and evidence grade. Nothing in this revision proves the planar
Jacobian conjecture. The exposition is expanded, in the step-by-step style of \S 5 of
Guccione--Guccione--Horruitiner--Valqui~\cite{GGHV}, so that each derivation shows its
intermediate equations rather than only its conclusion.
\end{center}

\begin{abstract}
Let $F = (P,Q) \colon \mathbb{A}^2 \to \mathbb{A}^2$ be a Keller map, that is,
$[P,Q] := P_x Q_y - P_y Q_x \in K^\times$ for $K$ algebraically closed of characteristic $0$,
and let $\lambda P + \mu Q$ be its pencil. The October 1 draft of this paper claimed that no
Keller map admits two distinct dicritical divisors, unconditionally. That claim is false: every
Keller map of degree at least $2$ has at least two exceptional dicriticals in total
(Corollary~\ref{cor:count-fails}), and the draft's $E_1$ degree obstruction
theorem is false as stated, with an explicit Keller
counterexample in \S\ref{sec:e1}.

The corrected target is \textbf{Theorem~R} (\S\ref{sec:statement}), which is \textbf{open}:
for every Keller map, the pencil $\lambda P + \mu Q$ has at most one dicritical divisor
\textbf{at infinity}. Theorem~R is equivalent to the planar Jacobian conjecture
(Proposition~\ref{prop:equiv}). This paper proves, or records as sketch or fixed-degree
computation, only the following steps toward it, each labelled with its grade:
\begin{itemize}[nosep]
\item the count of dicriticals (Proposition~\ref{prop:count}, grade B) and the position a
counterexample's dicriticals would occupy (\S\ref{sec:where});
\item the line at infinity is not dicritical for a non-affine Keller map
(Lemma~\ref{lem:linf}, grade B), and no first blow-up at infinity is dicritical
(Lemma~\ref{lem:e1}, algebra grade A, chart step grade B), so every dicritical at infinity
of a non-affine Keller map has depth at least $2$;
\item the $m = 1$ bridge equations collapse under two stated premises (P1) and (P2)
(Theorem~\ref{thm:bridge}, equations grade A, geometric reading grade B), where (P2) is
\textbf{not} derived and remains open;
\item the constant-leading-form argument for $m \geq 2$ on $E_k$, $k \geq 2$, as a
\textbf{proof sketch for free centers only} (Proposition~\ref{prop:constant-form}, grade C);
\item the adjacent $m = 1$ case as a computation at $d = 3$ only (grade C).
\end{itemize}
What remains open is collected in \S\ref{sec:open} (obligations O1--O6) and the tests any
proof must pass (T1--T4). \S\ref{sec:correspondence} maps each archived script to the exact
statement it certifies, under the corrected statements.
\end{abstract}

\vspace{1.2cm}
\hrule
\vspace{1.4cm}

\section{Introduction: setting, definitions, and the corrected statement}\label{sec:statement}

\subsection{Context}

The planar Jacobian conjecture asserts that a polynomial map
$(P,Q) \colon \mathbb{C}^2 \to \mathbb{C}^2$ with
$[P,Q] := P_x Q_y - P_y Q_x \in \mathbb{C}^\times$ is invertible. A \emph{Keller map} is
such a pair with constant nonzero Jacobian. This paper studies the dicritical divisors of the
pencil $\lambda P + \mu Q$ of a Keller map, and in particular how many of them can lie at
infinity. The October 1 draft claimed this line of work was closed. It was not, and this
revision says exactly what is proved, what is only sketched or computed at fixed degree, and
what is open.

Throughout, $K$ is an algebraically closed field of characteristic $0$; statements over any
field of characteristic $0$ reduce to its algebraic closure. We write
$F = (P,Q) \colon \mathbb{A}^2 \to \mathbb{A}^2$ for a Keller map, so
$[P,Q] = P_x Q_y - P_y Q_x \in K^\times$.

\subsection{Standard facts used}

We use the following standard facts, labelled (F1)--(F3) as in the restatement note:
\begin{itemize}[nosep]
\item[(F1)] $F$ is \'{e}tale. Its generic degree $e = [K(x,y) : K(P,Q)]$ is at least $1$.
\item[(F2)] The \emph{non-properness set} $S_F$ is the set of points of the target over which
$F$ is not proper. It is either empty or a curve whose components are rational
(Jelonek~\cite{Jelonek1993}, for dominant maps $\mathbb{A}^n \to \mathbb{A}^n$).
\item[(F3)] $F$ is an automorphism if and only if $S_F = \emptyset$. Indeed, if
$S_F = \emptyset$ then $F$ is a finite \'{e}tale cover of the simply connected
$\mathbb{A}^2$, hence an isomorphism; if $F$ is not an automorphism then $e \geq 2$, since
with $e = 1$ the map $F$ is birational and \'{e}tale, hence an open immersion by Zariski's
main theorem, so it is injective, and then bijective
(Bia\l{}ynicki-Birula--Rosenlicht~\cite{BBR1962}).
\end{itemize}

\subsection{Compactification, the pencil, and the standard definition}

Source and target are $\mathbb{A}^2 \subset \mathbb{P}^2$, with lines at infinity $L_\infty$
and $L'_\infty$. Let $\sigma \colon X \to \mathbb{P}^2$ be a sequence of point blow-ups over
$L_\infty$ such that $G = \bar{F} \circ \sigma \colon X \to \mathbb{P}^2$ is a morphism. The
\emph{curves at infinity} of $X$ are the components of $X \setminus \mathbb{A}^2$: the strict
transform of $L_\infty$ and the exceptional curves.

For $c \in \mathbb{A}^2$ (target), $\Lambda_c$ is the pencil whose members are $F^{-1}(\ell)$
for the lines $\ell \ni c$. Its members are the fibers of
$\varphi_c = \pi_c \circ F$, where $\pi_c$ is the projection from $c$. The paper's pencil
$\lambda P + \mu Q$ is $\Lambda_0$.

\begin{definition}[Dicritical divisor, standard definition]\label{def:dicritical}
A divisor $E$ over the source plane is \emph{dicritical} for $\Lambda_c$ if $\varphi_c$ is
not constant on $E$. Its \emph{degree} $m(E)$ is the degree of
$\varphi_c|_E \colon E \to \mathbb{P}^1$. The divisor $E$ is \emph{at infinity} if its center
on $\mathbb{P}^2$ lies on $L_\infty$; $L_\infty$ itself counts. Otherwise $E$ is
\emph{affine}. The \emph{depth} of $L_\infty$ is $0$; the exceptional curve of the blow-up at
a point $p$ has depth one more than the largest depth of the curves through $p$. Thus $E_1$
over a point of $L_\infty$ has depth $1$. A point of an exceptional curve is a
\emph{satellite} point if it lies on another component of the total transform of $L_\infty$,
and \emph{free} otherwise.
\end{definition}

\begin{remark}
The October 1 draft defined dicriticals as ``exceptional divisors on which the pulled-back
pencil has no base points.'' That is not the standard notion, and it is replaced by
Definition~\ref{def:dicritical} throughout this revision.
\end{remark}

\subsection{Counting dicriticals}

\begin{proposition}[The count]\label{prop:count}
\textit{Status: grade B, proof by hand from \textnormal{(F1)--(F3)}.}
Let $F$ be a Keller map and $c \in \mathbb{A}^2 \setminus S_F$.
\begin{enumerate}[label=(\alph*),nosep]
\item The affine dicriticals of $\Lambda_c$ are exactly the $e$ exceptional curves $E_q$ of
the blow-ups at the points $q \in F^{-1}(c)$. Each has degree $1$.
\item The dicriticals at infinity are the curves at infinity $D$ of $X$ for which $G(D)$ is a
curve. Then $G(D)$ is $L'_\infty$ or a component of $\bar{S}_F$, and
\begin{equation}\label{eq:count-degree}
m(D) = \deg\big(D \to G(D)\big) \cdot \deg G(D).
\end{equation}
\item $L'_\infty$ and every component of $\bar{S}_F$ are $G(D)$ for at least one such $D$.
So the number $N_\infty(c)$ of dicriticals at infinity satisfies
\begin{equation}\label{eq:count-lower}
N_\infty(c) \geq 1 + \#\{\text{components of } S_F\}.
\end{equation}
\end{enumerate}
\end{proposition}

\begin{proof}
We prove (a), (b), (c) in order, spelling out each step.

\emph{Part (a).} Since $c \notin S_F$, the map $F$ is proper over a neighbourhood of $c$.
By (F1), $F$ is \'{e}tale, so the fiber $F^{-1}(c)$ has exactly $e$ points, and $F$ is a local
isomorphism at each of them. Fix $q \in F^{-1}(c)$. Near $q$, the map $\varphi_c$ is the
projection from the point $c$ composed with a local isomorphism. One blow-up at $q$ resolves
it: in the blow-up chart, the exceptional curve $E_q$ parametrizes the directions through $q$,
and $\varphi_c$ restricts to the projection itself, hence maps $E_q$ isomorphically onto
$\mathbb{P}^1$. So $E_q$ is dicritical of degree $1$. At every other affine point,
$\varphi_c$ is regular, so no other affine divisor is dicritical. This gives exactly the $e$
exceptional curves $E_q$, each of degree $1$.

\emph{Part (b).} Let $D$ be a curve at infinity. First, $G(D)$ is not the point $c$: the map $G$
sends no curve at infinity into $\mathbb{A}^2 \setminus S_F$, and $c \notin S_F$ by hypothesis.
If $G(D)$ is a point different from $c$, then $\varphi_c = \pi_c \circ G$ is constant on $D$,
since $\pi_c$ sends every point other than $c$ to a single point of $\mathbb{P}^1$ only after
fixing the line through $c$ --- more precisely, $\pi_c$ is constant on no curve through a point
other than $c$ except by collapsing it, and here $G(D)$ is a single point $\neq c$, so
$\varphi_c|_D$ is that constant value. Hence $D$ is not dicritical in this case.

If instead $G(D)$ is a curve, then either $G(D) = L'_\infty$, or $G(D)$ meets
$\mathbb{A}^2$. In the second case, every point of $G(D) \cap \mathbb{A}^2$ is a limit of $F$
along a path to infinity, so it lies in $S_F$; since $G(D)$ is irreducible, $G(D)$ is a
component of the closure $\bar{S}_F$. Neither $L'_\infty$ nor a component of $S_F$ is a line
through $c$. Therefore the projection $\pi_c$ restricted to the curve $G(D)$ has degree
$\deg G(D)$, and composing with $D \to G(D)$ gives
\[
m(D) = \deg(\varphi_c|_D) = \deg\big(D \to G(D)\big) \cdot \deg G(D),
\]
which is \eqref{eq:count-degree}. So $D$ is dicritical exactly when $G(D)$ is a curve, with
the stated degree.

\emph{Part (c).} The map $G$ is proper and surjective, so $L'_\infty$ is dominated by some
curve at infinity: there is $D$ with $G(D) = L'_\infty$. By the description of $S_F$ as
$G(X \setminus \mathbb{A}^2) \cap \mathbb{A}^2$, each component of $S_F$ is likewise
dominated by a curve at infinity. Counting the distinct images $L'_\infty$ and the components
of $S_F$ gives the lower bound \eqref{eq:count-lower}.
\end{proof}

\begin{corollary}[Automorphisms]\label{cor:automorphisms}
\textit{Status: grade B.}
If $F$ is an automorphism, then $\Lambda_c$ has exactly two dicriticals for every $c$: one
affine, over $F^{-1}(c)$, and one at infinity, the unique curve $D$ with
$G(D) = L'_\infty$ (here $G$ is birational). Both have degree $1$.
\end{corollary}

\begin{proof}
If $F$ is an automorphism, then $e = 1$ and $S_F = \emptyset$ by (F3). Part (a) of
Proposition~\ref{prop:count} gives exactly one affine dicritical, over the single point
$F^{-1}(c)$, of degree $1$. Part (c) with $S_F = \emptyset$ gives at least one dicritical at
infinity, with image $L'_\infty$; since $G$ is birational, there is exactly one such curve,
and \eqref{eq:count-degree} gives its degree as $1 \cdot 1 = 1$.
\end{proof}

\subsection{The October 1 main theorem is false}

\begin{corollary}[The theorem as worded fails]\label{cor:count-fails}
\textit{Status: grade B in general; exact in the examples below.}
Every Keller map of degree at least $2$ has at least two exceptional dicriticals in
$\Lambda_0$, after a translation of the target that puts $0$ in
$F(\mathbb{A}^2) \setminus S_F$. An automorphism needs no translation. Hence the October 1
draft's theorem (``no two distinct dicritical divisors,'' with no at-infinity restriction) is
false for every Keller map of degree at least $2$. If only exceptional curves count as
divisors of the resolution, it holds exactly for affine maps, whose dicritical at infinity is
$L_\infty$ itself; if $L_\infty$ counts, it fails for them too.
\end{corollary}

\begin{proof}
Let $F$ have degree at least $2$. Translate the target so that $0 \notin S_F$ and
$0 \in F(\mathbb{A}^2)$; then $F^{-1}(0)$ is nonempty. By Proposition~\ref{prop:count}(a),
each point of $F^{-1}(0)$ contributes an affine dicritical of degree $1$, and by part (c)
there is at least one dicritical at infinity. That is already two dicriticals in total. The
example below exhibits both of them explicitly for an automorphism, where no translation is
needed.
\end{proof}

\begin{example}[Two dicriticals of an automorphism]\label{ex:automorphism}
Take
\begin{equation}\label{eq:auto-map}
P = x - y^2, \qquad Q = y.
\end{equation}
We compute the Jacobian step by step. First the partial derivatives:
\[
P_x = 1, \qquad P_y = -2y, \qquad Q_x = 0, \qquad Q_y = 1.
\]
Hence
\[
[P,Q] = P_x Q_y - P_y Q_x = 1 \cdot 1 - (-2y) \cdot 0 = 1 \in K^\times,
\]
so $F = (P,Q)$ is a Keller map. It is an automorphism: solving $P = x - y^2$ and $Q = y$
gives $y = Q$ and $x = P + Q^2$, so the inverse is $(P,Q) \mapsto (P + Q^2, Q)$, and indeed
$P + Q^2 = (x - y^2) + y^2 = x$.

\emph{The affine dicritical.} The fiber over $0$ is the single point $(0,0)$, since
$Q = y = 0$ and $P = x - y^2 = 0$ force $x = 0$. Blow up $\mathbb{A}^2$ at the origin and
use the chart
\[
x = s, \qquad y = s t,
\]
so the exceptional divisor $E$ is $\{s = 0\}$ with coordinate $t$. Substituting:
\begin{align}
P &= s - (st)^2 = s - s^2 t^2 = s\,(1 - s t^2), \\
Q &= s t.
\end{align}
The map $\varphi_0$ is the projection from $0$, so on this chart it is the ratio
\[
\varphi_0 = [P : Q] = [\,s(1 - s t^2) : s t\,] = [\,1 - s t^2 : t\,].
\]
Restricting to the exceptional divisor, where $s = 0$, gives
\[
\varphi_0|_E = [1 : t],
\]
which is non-constant in $t$. Hence $E$ is dicritical, of degree $1$.

\emph{The dicritical at infinity.} By Corollary~\ref{cor:automorphisms}, $\Lambda_0$ has
exactly one dicritical at infinity, the unique curve $D$ with $G(D) = L'_\infty$, also of
degree $1$. Together with the affine dicritical above, this is two distinct dicriticals.
The October 1 theorem, which forbade two distinct dicritical divisors without any
at-infinity restriction, is therefore false, already for automorphisms.
\end{example}

\subsection{The corrected statement: Theorem R}

\begin{theorem}[Theorem R --- the corrected main theorem; OPEN]\label{thm:R}
\textit{Status: open, not proved in this paper or elsewhere.}
For every Keller map $(P,Q)$, the pencil $\lambda P + \mu Q$ has at most one dicritical
divisor at infinity.
\end{theorem}

\begin{proposition}\label{prop:onedic}
\textit{Status: grade B.}
Let $F$ be a Keller map and $c \notin S_F$. Then $\Lambda_c$ has at most one dicritical at
infinity if and only if $F$ is an automorphism. In that case it has exactly one, of degree $1$.
\end{proposition}

\begin{proof}
If $F$ is an automorphism, Corollary~\ref{cor:automorphisms} gives exactly one dicritical at
infinity, of degree $1$. Conversely, suppose $\Lambda_c$ has at most one dicritical at
infinity. By Proposition~\ref{prop:count}(c), the number of dicriticals at infinity is at
least $1 + \#\{\text{components of } S_F\}$. At most one therefore forces
$\#\{\text{components of } S_F\} = 0$, that is, $S_F = \emptyset$. By (F3), $F$ is an
automorphism.
\end{proof}

\begin{proposition}[Theorem R is equivalent to the planar Jacobian conjecture]\label{prop:equiv}
\textit{Status: grade B.}
Theorem~R holds for all Keller maps if and only if the planar Jacobian conjecture holds.
\end{proposition}

\begin{proof}
If the conjecture holds, then every Keller map is an automorphism, and
Corollary~\ref{cor:automorphisms} gives exactly one dicritical at infinity, in particular at
most one. So Theorem~R holds.

Conversely, let $F$ be a Keller map and pick $c \notin S_F$. Such a $c$ exists because $S_F$
is either empty or a curve, never the whole target. Then $F - c$ is a Keller map with
$0 \notin S_{F-c}$, and its pencil $\lambda (P - c_1) + \mu (Q - c_2)$ is $\Lambda_c(F)$.
Applying Theorem~R to $F - c$ and Proposition~\ref{prop:onedic} to $\Lambda_c(F)$ makes $F$
an automorphism. Since $F$ was arbitrary, the conjecture holds.
\end{proof}

\begin{remark}
Proposition~\ref{prop:equiv} sets the bar for this paper: a complete proof of the corrected
theorem would be a proof of the conjecture. Everything that follows is therefore measured
against that bar, and no step below is presented as closing it.
\end{remark}

\subsection{Where a counterexample's dicriticals would sit}\label{sec:where}

Let $F$ be a Keller map that is not an automorphism, and let $c \notin S_F$. Then
$\Lambda_c$ has:
\begin{itemize}[nosep]
\item $e \geq 2$ affine dicriticals, by (F3) and Proposition~\ref{prop:count}(a);
\item at least one $E_\infty$ with $G(E_\infty) = L'_\infty$, of degree
$\deg(E_\infty \to L'_\infty)$;
\item at least one $E_S$ over each component $C$ of $S_F$, of degree
$\deg(E_S \to C) \cdot \deg C$, by \eqref{eq:count-degree}.
\end{itemize}
$E_\infty$ always exists, so Theorem~R amounts to showing that $E_S$ cannot exist. The two
halves of the October 1 draft, read against this picture, were: the first half claimed there
is no dicritical of degree $m \geq 2$ for any Keller map, which applied to $E_S$ would
already exclude every counterexample whose $S_F$ has a component of degree at least $2$; the
second half claimed there are no two dicriticals of degree $1$, which would handle the rest.
Either half, to succeed, must apply to $E_\infty$ and $E_S$ wherever they sit in the
resolution tree: over any point of $L_\infty$, at any depth, and after free or satellite
blow-ups. Sections~\ref{sec:depth01}--\ref{sec:bridge} record how far each half actually goes.

\section{Dicriticals at infinity of small depth}\label{sec:depth01}

This section proves the two depth $0$ and depth $1$ steps toward Theorem~R, each with every
intermediate computation shown, and then records the constant-leading-form argument in its
honest state: a sketch for free centers only. The conclusion of this section is that every
dicritical at infinity of a non-affine Keller map has depth at least $2$. Nothing here touches
depth $2$ or beyond, except for the sketch in \S\ref{sec:constant-form}, whose limits are
stated there.

\subsection{The line at infinity}\label{sec:linf}

\begin{lemma}[The line at infinity]\label{lem:linf}
\textit{Status: grade B, proof by hand; the binary-form step is executed in
\filename{base\_point\_at\_infinity\_check.py}, check $1$.}
For a non-affine Keller map, $L_\infty$ is not dicritical for any $\Lambda_c$.
\end{lemma}

\begin{proof}
Let $d = \max(\deg P, \deg Q) \geq 2$, the non-affine case. Write $P_d$ and $Q_d$ for the
homogeneous components of degree $d$ of $P$ and $Q$ (zero if the corresponding degree is
smaller than $d$). The restriction $\varphi_c|_{L_\infty}$ is the rational function
$P_d / Q_d$ on $L_\infty \cong \mathbb{P}^1$, where the point $c$ contributes only a constant
shift that does not affect constancy. We show this ratio is constant, in two cases.

\emph{Case 1: unequal degrees.} Suppose $\deg P \neq \deg Q$, say $\deg P > \deg Q$ (the other
case is symmetric). Then on $L_\infty$, the term of highest degree in the pencil is $P_d$
alone, so $P_d / Q_d$ is $\infty$ wherever $Q_d$ vanishes to lower order, and the restriction
is the constant $\infty$. Explicitly: homogenizing, the leading part of a member
$\lambda P + \mu Q$ at infinity is $\lambda P_d$ (since $\deg Q < d$), so the pencil has the
single point at infinity determined by $P_d$ as its value, independent of the member. The
restriction is constant, so $L_\infty$ is not dicritical. The case $\deg Q > \deg P$ gives
the constant $0$ by the same computation.

\emph{Case 2: equal degrees.} Suppose $\deg P = \deg Q = d$. The Jacobian $[P,Q]$ is a nonzero
constant, so in particular its homogeneous component of degree $2d - 2$ vanishes. That
component is exactly $[P_d, Q_d]$, because differentiating lowers degree by $1$ in each factor
and the bracket of the top components is the top component of the bracket:
\begin{equation}\label{eq:top-bracket}
[P,Q]_{2d-2} = [P_d, Q_d] = 0.
\end{equation}
So it suffices to prove the following binary-form lemma.

\begin{lemma}[Proportionality of binary forms with vanishing bracket]\label{lem:binary}
Let $f, g$ be binary forms (homogeneous polynomials in two variables) of the same degree
$\geq 1$. If $[f,g] = 0$, then $f = r g$ for some constant $r \in K$.
\end{lemma}

\begin{proof}[Proof of Lemma~\ref{lem:binary}, step by step]
Work in the rational function field and set $h = f / g$, a rational function of degree $0$,
meaning $h(tx, ty) = h(x,y)$ for all $t$. Differentiating this homogeneity relation with
respect to $t$ at $t = 1$ gives Euler's relation for $h$:
\begin{equation}\label{eq:euler-h}
x h_x + y h_y = 0.
\end{equation}
Next, the bracket is bilinear and satisfies the quotient rule in the first argument:
\[
[h, g] = \frac{[f,g]}{g} - \frac{f [g,g]}{g^2} = \frac{[f,g]}{g} = 0,
\]
using $[g,g] = 0$ and the hypothesis $[f,g] = 0$. Writing the bracket out,
\begin{equation}\label{eq:bracket-hg}
h_x g_y - h_y g_x = 0.
\end{equation}
Now suppose $(h_x, h_y) \neq (0,0)$. Equation~\eqref{eq:bracket-hg} says the vector
$(h_x, h_y)$ is proportional to $(g_x, g_y)$: both are orthogonal to the same direction, so
there is a scalar $\lambda$ with $(g_x, g_y) = \lambda (h_x, h_y)$ wherever the vectors are
defined. Equation~\eqref{eq:euler-h} says $(h_x, h_y)$ is proportional to $(-y, x)$: indeed,
$x h_x + y h_y = 0$ means $(h_x, h_y)$ is orthogonal to $(x,y)$, and in two variables the
orthogonal direction is spanned by $(-y, x)$. Combining the two proportionalities,
$(g_x, g_y)$ is proportional to $(-y, x)$, so
\[
x g_x + y g_y = 0.
\]
But Euler for the homogeneous form $g$ of degree $\deg g \geq 1$ says
\[
x g_x + y g_y = (\deg g) \cdot g \neq 0.
\]
This is a contradiction. Hence the supposition is false, and $(h_x, h_y) = (0,0)$. In
characteristic $0$, a rational function with vanishing partial derivatives is constant, so
$h = r$ for some $r \in K$, that is, $f = r g$.
\end{proof}

Returning to Case 2: Lemma~\ref{lem:binary} applied to $f = P_d$, $g = Q_d$ gives
$P_d = r Q_d$, so $\varphi_c|_{L_\infty} = P_d / Q_d = r$ is constant. Hence $L_\infty$ is not
dicritical. This completes both cases.
\end{proof}

\subsection{First blow-ups at infinity}\label{sec:e1lemma}

\begin{lemma}[First blow-ups at infinity]\label{lem:e1}
\textit{Status: algebra grade A, in \filename{lean/TwoDicritical/E1Keller.lean}
(\texttt{e1\_lead\_proportional}, \texttt{jac2\_eq\_zero\_of\_natDegree\_zero}); chart step
grade B, executed in \filename{lean/check\_e1\_index.py}, check $5$, not Lean-formalized.}
For a Keller map, the exceptional curve of the blow-up at any point of $L_\infty$ is not
dicritical. This holds for every degree $m$, not only for $m = 2$.
\end{lemma}

\begin{proof}
A linear change of coordinates keeps the Keller condition (the Jacobian transforms by the
constant determinant of the change), so we may move the point of $L_\infty$ to
$[1 : 0 : 0]$. Let $E_1$ be the exceptional curve of the blow-up at this point, with the
chart coordinate $t_1$ on $E_1$ as in \filename{lean/check\_e1\_index.py}. The proof has three
steps: the restriction formula, the Wronskian computation, and the degenerate cases.

\emph{Step 1: the restriction formula.} Suppose $P$ and $Q$ have the same $x$-degree
$n \geq 1$, with $x^n$-coefficients $A, D \in K[y]$, so that
\begin{equation}\label{eq:lead-expansion}
P = A(y)\, x^n + (\text{terms of $x$-degree} < n), \qquad
Q = D(y)\, x^n + (\text{terms of $x$-degree} < n).
\end{equation}
In the $E_1$ chart, the ratio $P/Q$ restricts to $E_1$ as
\begin{equation}\label{eq:e1-restriction}
\frac{P}{Q}\bigg|_{E_1} = \frac{A(1/t_1)}{D(1/t_1)}.
\end{equation}
The reason is that on $E_1$ only the top $x$-degree terms survive: lower $x$-degree terms
vanish to higher order along $E_1$, and the chart coordinate $t_1$ inverts the $y$-variable,
which is why $A$ and $D$ are evaluated at $1/t_1$. If instead the $x$-degrees of $P$ and $Q$
differ, say $\deg_x P > \deg_x Q$, then the numerator dominates and the restriction is the
constant $\infty$; if $\deg_x P < \deg_x Q$, it is the constant $0$. In both differing cases
$E_1$ is already not dicritical, so only the equal-degree case remains.

\emph{Step 2: the Wronskian computation.} Differentiate the expansions
\eqref{eq:lead-expansion} term by term. The top $x$-degree contributions are:
\begin{align}
P_x &= n A(y)\, x^{n-1} + (\text{$x$-degree} < n-1), \\
P_y &= A'(y)\, x^{n} + (\text{$x$-degree} < n), \\
Q_x &= n D(y)\, x^{n-1} + (\text{$x$-degree} < n-1), \\
Q_y &= D'(y)\, x^{n} + (\text{$x$-degree} < n).
\end{align}
Now form the Jacobian and collect the top $x$-degree. The product $P_x Q_y$ has top
$x$-degree $(n-1) + n = 2n - 1$, with coefficient
\[
(n A) \cdot D' = n A D'.
\]
The product $P_y Q_x$ likewise has top $x$-degree $n + (n-1) = 2n - 1$, with coefficient
\[
A' \cdot (n D) = n A' D.
\]
Every other product of a term from $P_x, P_y$ with a term from $Q_x, Q_y$ has $x$-degree
strictly less than $2n - 1$: lowering the $x$-degree of either factor by at least $1$ lowers
the total by at least $1$. Therefore the $x^{2n-1}$-coefficient of $[P,Q]$ is exactly
\begin{equation}\label{eq:wronskian-coeff}
[P,Q]_{x^{2n-1}} = n\,\big(A D' - A' D\big).
\end{equation}
Since $(P,Q)$ is Keller, $[P,Q]$ is a nonzero constant, so every coefficient of a positive
power of $x$ in $[P,Q]$ vanishes. In particular the coefficient \eqref{eq:wronskian-coeff}
vanishes, and since $n \geq 1$,
\[
A D' - A' D = 0.
\]
This is the vanishing of the Wronskian of $A$ and $D$. In characteristic $0$ it forces $A$
and $D$ to be proportional: indeed, where $D \neq 0$,
\[
\left(\frac{A}{D}\right)' = \frac{A' D - A D'}{D^2} = 0,
\]
so $A/D$ is constant, say $A = r D$ for $r \in K$. (This is the Lean lemma
\texttt{e1\_lead\_proportional}.) Substituting into the restriction formula
\eqref{eq:e1-restriction} gives
\[
\frac{P}{Q}\bigg|_{E_1} = \frac{r D(1/t_1)}{D(1/t_1)} = r,
\]
a constant. Hence $E_1$ is not dicritical in the equal-degree case either.

\emph{Step 3: the degenerate case.} If both $x$-degrees are $0$, then $P, Q \in K[y]$ are
functions of $y$ alone, so $P_x = Q_x = 0$ and $[P,Q] = 0$. This contradicts the Keller
condition $[P,Q] \in K^\times$ (Lean lemma \texttt{jac2\_eq\_zero\_of\_natDegree\_zero}). So
this case cannot occur for a Keller map.

Since a linear change of coordinates moves any point of $L_\infty$ to $[1:0:0]$ while
preserving the Keller condition, the same conclusion holds over every point at infinity.
\end{proof}

\begin{corollary}[Depth at least $2$]\label{cor:depth2}
\textit{Status: follows from Lemmas~\ref{lem:linf} and~\ref{lem:e1}.}
Every dicritical at infinity of a non-affine Keller map has depth at least $2$.
\end{corollary}

\begin{proof}
Depth $0$ is $L_\infty$ itself, excluded by Lemma~\ref{lem:linf} for non-affine maps.
Depth $1$ curves are the exceptional curves of blow-ups at points of $L_\infty$, excluded by
Lemma~\ref{lem:e1}. The smallest remaining depth is $2$.
\end{proof}

\subsection{The constant-leading-form argument: a sketch for free centers only}\label{sec:constant-form}

\begin{proposition}[Constant leading form --- proof sketch]\label{prop:constant-form}
\textit{Status: grade C, proof sketch; free centers only. This is not a proof for satellite
chains or for depth $\geq 2$ in general (obligation O3).}
Suppose $E_k$ ($k \geq 2$) is dicritical of order $m \geq 1$ for the pencil. Then after the
standard normalization, the leading coefficient satisfies $\deg_{u'} U(u',0) = 0$.
\end{proposition}

\begin{proof}[Proof sketch]
Let $E_k$ ($k \geq 2$) be an exceptional divisor in the resolution, with local coordinates
$(u',v')$ in which $E_k = \{v' = 0\}$. For a polynomial $S \in K[x,y]$ of degree $d$, write
the pullback as
\begin{equation}\label{eq:pullback}
S = {v'}^{-\nu}\,U(u',v')\cdot(\text{monomial in }u'),
\end{equation}
where $\nu = \nu_{E_k}(S)$ is the order of vanishing and $U$ is a unit (nonzero at the
center). The dicritical condition of order $m$ requires the $u'$-degree of the pencil
restriction to equal $m$.

Track the unit $U$ through the blow-up sequence, one blow-up at a time. Consider a single
blow-up at a center with $u_0 \neq 0$ (a \emph{free} center: the center does not lie on another
component of the total transform, so the local coordinate $u'$ is nonzero there). In the
blow-up chart, substituting the center shifts the leading coefficient by the constant factor
$u_0^{-2}$: schematically, if before the blow-up the leading term is $c \cdot {u'}^j$, then
after the blow-up it is $c\, u_0^{-2} \cdot {u''}^{j}$ in the new coordinate $u''$, with the
same exponent $j$ and a nonzero constant multiple. The point is that multiplication by the
nonzero constant $u_0^{-2}$ changes no $u'$-degree: the exponent $j$ is carried through
unchanged, and no new $u'$-dependence is created, because the center coordinate $u_0$ is a
nonzero constant rather than a function of $u'$.

Repeating this step along the chain: each blow-up at a center with $u_0 \neq 0$ multiplies
the leading coefficient by a nonzero constant and preserves its $u'$-degree. A direct
computation in the normalized chart then shows the $u'$-degree of $U(u',0)$ is $0$. The
verification of the degree count is in \filename{general\_constant\_form.py}
(\S\ref{sec:correspondence}).

\emph{Where the sketch stops.} The step above uses $u_0 \neq 0$ at every blow-up. If a center
lies on another component of the total transform (a satellite center), then the corresponding
coordinate vanishes at the center, the factor is no longer a nonzero constant, and the
degree-tracking argument as written does not go through. The sketch therefore covers free
chains only. It does not prove that no $E_k$ ($k \geq 2$) is dicritical of order
$m \geq 2$ in general, and it is not presented as doing so: the general $m \geq 2$ statement
at depth $\geq 2$ is obligation O3 in \S\ref{sec:open}. The middle-divisor ($E_2$) and leaf
cases ($E_3$, $E_4$) are checked separately as computations in \filename{middle\_e2.py},
\filename{leaf\_e3.py}, and \filename{leaf\_e4.py}; those are fixed instances, not the
general argument.
\end{proof}

\begin{remark}
In the October 1 draft, Proposition~\ref{prop:constant-form} was followed by a corollary
claiming no $E_k$ ($k \geq 2$) is dicritical of order $m \geq 2$ for any Keller map, and that
corollary was then used to close the $m \geq 2$ half of the main theorem. That inference is
withdrawn: a sketch for free centers does not yield the general corollary, and the general
corollary is not claimed anywhere in this revision.
\end{remark}

\section{The first exceptional divisor: a false obstruction and its correction}\label{sec:e1}

This section does two things the October 1 draft got backwards. First, it exhibits the draft's
$E_1$ degree obstruction as false, shows the exact step where its proof goes wrong, and
computes the counterexample in full. Second, it records the corrected statement, which was
already proved as Lemma~\ref{lem:e1} in \S\ref{sec:e1lemma}: first blow-ups at infinity are
never dicritical for a Keller map, for every degree $m$. The correction is strictly stronger
than the draft's claim in scope (every $m$, not only $m = 2$) and strictly weaker in what it
asserts about $E_1$ itself (it does not claim anything about a hypothetical $m = 2$ open
subset, which, as the counterexample shows, contains Keller pairs).

\subsection{What the October 1 draft claimed}

The draft's $E_1$ degree obstruction theorem read, in substance: let
$P, Q$ lie in the $E_1$ $m = 2$ open subset, meaning $\nu_{E_1} = -2$ and both pencil
restrictions have exact degree $2$; then $[P,Q]$ is never a nonzero constant. Its proof argued
as follows: every monomial $x^J y^K$ allowed by the order condition $\nu_{E_1} \geq -2$
satisfies $J + K \geq 2$; hence $P_x, P_y, Q_x, Q_y$ have no constant term; hence $[P,Q]$ has
minimum total degree $2$ and cannot be a nonzero constant. Every step after the first
sentence is correct arithmetic. The first sentence is false, and this subsection shows why,
step by step.

\subsection{The order condition, computed correctly}

\begin{lemma}[The $E_1$ order condition]\label{lem:e1-order}
In the $E_1$ chart, the valuation satisfies
\begin{equation}\label{eq:e1-valuation}
\nu_{E_1}(x^i y^j) = -i.
\end{equation}
Consequently $\nu_{E_1} \geq -2$ holds for a monomial $x^i y^j$ if and only if $i \leq 2$,
that is, if and only if the $x$-degree is at most $2$. It says nothing about the total degree
$i + j$.
\end{lemma}

\begin{proof}
Equation~\eqref{eq:e1-valuation} is the normalization of the $E_1$ valuation used in the
$E_1$ chart (see \filename{lean/check\_e1\_index.py} and
\filename{lean/TwoDicritical/E1Keller.lean}): the variable $x$ has a pole of order $1$ along
$E_1$, contributing $-1$ per power of $x$, while $y$ contributes $0$. Hence for the monomial
$x^i y^j$,
\[
\nu_{E_1}(x^i y^j) = i \cdot \nu_{E_1}(x) + j \cdot \nu_{E_1}(y) = i \cdot (-1) + j \cdot 0 = -i.
\]
The inequality $\nu_{E_1} \geq -2$ is therefore
\[
-i \geq -2 \quad\Longleftrightarrow\quad i \leq 2.
\]
In words: the condition bounds the $x$-degree by $2$. It does not bound $j$, and it does not
bound the total degree $i + j$ from below: the monomial $y$ itself, with $i = 0$ and $j = 1$,
satisfies $\nu_{E_1}(y) = 0 \geq -2$ while having total degree $1$. The draft's claim that
every allowed monomial has $J + K \geq 2$ is thus false on its face, once the valuation is
written out.
\end{proof}

\subsection{How the mix-up happened}

The draft's error is an index mix-up, identified in the Lean analysis
(\filename{lean/README.md}, item 2). The computation \filename{e1\_solitary.py} indexes its
unknowns $p[J,K]$ by the homogeneous monomial $X^{d-J-K} Y^J Z^K$, that is, by the affine
monomial
\[
x^{d-J-K} y^J.
\]
The draft's proof, and the scripts \filename{e1\_keller\_search.py} and the Lean file
\filename{lean/TwoDicritical/E1Obstruction.lean}, read $p[J,K]$ as the monomial $x^J y^K$
instead. Under that misreading, the condition on the indices becomes a condition on the total
degree $J + K$, and the draft concluded that every allowed monomial has total degree at least
$2$. At $d = 4$, both index sets have $12$ monomials, so the count agreed and the mix-up
was hidden: the $12$-dimensional span at $d = 4$ is real, but it is a span of monomials of
$x$-degree at most $2$, not of total degree at least $2$. The file
\filename{lean/TwoDicritical/E1Obstruction.lean} proves a true lemma about the misread set of
monomials (if every monomial of $P$ and $Q$ has total degree at least $2$, then $[P,Q]$ has
no constant term), but that set is not the $E_1$ condition, so the lemma does not prove
anything about $E_1$.

\subsection{The counterexample, computed in full}

\begin{example}[A Keller pair in the draft's $E_1$ open subset]\label{ex:e1-counter}
\textit{Status: grade A in Lean (\filename{lean/TwoDicritical/E1Keller.lean},
\texttt{thm\_e1\_false}); the computation below is exact.}
Set
\begin{equation}\label{eq:e1-counter}
P = (x + y^2)^2 + y, \qquad Q = P + x + y^2.
\end{equation}
We expand $P$ first, term by term:
\[
P = x^2 + 2 x y^2 + y^4 + y.
\]
Its partial derivatives are:
\begin{align}
P_x &= 2x + 2y^2 = 2\,(x + y^2), \\
P_y &= 4x y + 4y^3 + 1 = 4y\,(x + y^2) + 1.
\end{align}
Now compute $[P,Q]$. Since $Q = P + x + y^2$ and $[P,P] = 0$, bilinearity of the bracket
gives
\[
[P,Q] = [P, x] + [P, y^2].
\]
Each bracket is computed separately. For $[P,x]$: $x$ has derivatives $1$ and $0$, so
\[
[P,x] = P_x \cdot 0 - P_y \cdot 1 = -P_y = -\big(4y(x+y^2) + 1\big).
\]
For $[P, y^2]$: $y^2$ has derivatives $0$ and $2y$, so
\[
[P, y^2] = P_x \cdot 2y - P_y \cdot 0 = 2y P_x = 2y \cdot 2(x+y^2) = 4y(x+y^2).
\]
Adding the two contributions, the $4y(x+y^2)$ terms cancel exactly:
\begin{equation}\label{eq:e1-bracket}
[P,Q] = -\big(4y(x+y^2) + 1\big) + 4y(x+y^2) = -1 \in K^\times.
\end{equation}
So $(P,Q)$ is a Keller pair.

It lies in the draft's $E_1$ $m = 2$ open subset at $d = 4$: its $x$-degree is $2$ (the term
$x^2$), so $\nu_{E_1} = -2$ by Lemma~\ref{lem:e1-order}, and each pencil restriction is the
function $C(t_1) = t_1^2$ of exact degree $2$ (this is the membership statement formalized as
\texttt{thm\_e1\_false} in \filename{lean/TwoDicritical/E1Keller.lean}, with the open subset
denoted \texttt{E1Open4} there). Yet $[P,Q] = -1$ is a nonzero constant, contradicting the
draft's theorem directly.

Two further observations pin down the failure. First, $P$ contains the monomial $y$, of total
degree $1$: the draft's claimed lower bound $J + K \geq 2$ fails for a monomial that is
actually present. Second, $P/Q$ restricts to $1$ on $E_1$ for this pair, so $E_1$ is not
dicritical for them at all: the draft's ``open subset'' does not even capture an $m = 2$
dicritical. The draft's proof argued about the wrong monomials, and its conclusion is false.
\end{example}

\begin{remark}[Withdrawn theorem]
The draft's $E_1$ degree obstruction theorem and its corollary (``no dicritical of order
$m \geq 2$ exists for any Keller map'') are both withdrawn. The corollary depended on the false theorem for the $E_1$ case and on the
free-only sketch of Proposition~\ref{prop:constant-form} for the $E_k$, $k \geq 2$, case;
neither leg supports it as stated.
\end{remark}

\subsection{The corrected \texorpdfstring{$E_1$}{E1} statement}

The intended conclusion of the draft's $E_1$ argument --- that the first blow-up cannot be
dicritical for a Keller map --- is true, and it is Lemma~\ref{lem:e1} of
\S\ref{sec:e1lemma}, proved there for every degree $m$. To recap the logic in one place: the
restriction of $P/Q$ to $E_1$ is $A(1/t_1)/D(1/t_1)$ when the $x$-degrees agree (and $0$ or
$\infty$ when they differ), and the Keller condition forces the Wronskian $AD' - A'D$ to
vanish via the $x^{2n-1}$-coefficient \eqref{eq:wronskian-coeff}, hence $A = rD$ and the
restriction is the constant $r$. The algebra is Lean-grade-A; the chart step is executed, not
Lean-formalized (obligation O5). No $m = 2$ open subset, and no total-degree argument, is
used anywhere in the corrected proof.

\section{The \texorpdfstring{$m=1$}{m=1} bridge, under its premises}\label{sec:bridge}

This section keeps the October 1 draft's Vandermonde and block-rank argument, which is correct
as linear algebra, but restates its scope exactly. The argument proves a collapse of a system
of equations. Whether that system is the system of two $m = 1$ dicriticals depends on two
premises, (P1) and (P2), stated below. Premise (P2) is not derived anywhere in this paper; it
is an open obligation (O2). The theorem of this section is therefore a theorem about equations
under premises, not an unconditional geometric statement, and its corollaries are fenced
accordingly.

\subsection{The chart and the two premises}

Read through its chart, the bridge concerns the following geometry. The paper's pullback is
\[
S^h\big(1, b, b^2(a + \tilde{a} b)\big),
\]
the chart at the point $w_2 = a$ of $E_2$ over $[1 : 0 : 0]$, where $w_2 = x / y^2$. The two
centers $a_3 \neq a_4$ are points of this $E_2$. The bridge argument rests on:
\begin{itemize}[nosep]
\item[(P1)] the two dicriticals lie over two distinct \emph{free} points of this $E_2$;
\item[(P2)] every member of the pencil has pole order at most $1$ on them.
\end{itemize}
Premise (P1) is part of the geometric reading of the equations: the chart above is the chart
at a free point of $E_2$, and the argument is set in that chart. Premise (P2) is a hypothesis
on the members, not a consequence of the definitions used here: the paper defines $m$ as the
degree of $\varphi_c|_E$, and degree $1$ does not force pole order $1$
(\filename{lean/check\_satellite.py}, check $7$). Deriving (P2) from the definition of $m$,
or supplying a separate argument for higher pole order, is obligation O2 in
\S\ref{sec:open}. Nothing in this section assumes (P2) silently: every statement that needs
it says so.

\subsection{Setup of the equations}

Consider two distinct centers $a_3 \neq a_4$ (both nonzero), each requiring order
$\operatorname{ord}_b(S) \geq 2d - 1$ (the $m = 1$ condition, under (P2)). For
$S = \sum p_{JK} x^J y^K$, the pullback at center $a$ is
\begin{equation}\label{eq:bridge-pullback}
S_a(b) = \sum_{J,K} p_{JK}\,(a + \tilde{a} b)^{d-J-K}\, b^{2d-2J-K},
\end{equation}
and we require the coefficients of $b^0, \dots, b^{2d-2}$ to vanish at $a = a_3$ and at
$a = a_4$. The exponent $d - J - K$ is the remaining degree after the monomial $x^J y^K$ is
factored out, and the exponent $2d - 2J - K$ is the $b$-order contributed by that monomial in
the chart; the condition $\operatorname{ord}_b(S) \geq 2d - 1$ is exactly the vanishing of
all coefficients below $b^{2d-1}$.

\begin{theorem}[Bridge collapse, under (P1) and (P2)]\label{thm:bridge}
\textit{Status: equations grade A, in \filename{lean/TwoDicritical/BridgeCollapse.lean};
geometric reading grade B, executed in \filename{lean/check\_satellite.py}.}
For arbitrary distinct nonzero $a_3, a_4$ and arbitrary $d \geq 2$, the simultaneous
conditions $\operatorname{ord}_b(S) \geq 2d - 1$ at both centers force $p_{JK} = 0$ whenever
$J \geq 1$. The nullspace is exactly $\operatorname{span}\{1, y\}$.
\end{theorem}

In particular, every $P, Q$ in the span is a function of $y$ alone, so $[P,Q] \equiv 0$: no
Keller pair exists in the system of equations. Whether this rules out two geometric $m = 1$
dicriticals depends on (P1) and (P2); see \S\ref{sec:bridge-scope}.

\subsection{Descending induction on weighted degree}

Order monomials by the weighted degree $W = 2J + K$, descending. For fixed $W$, write the
variables as $(J_k, K_k)$, $k = 1, \dots, n$, ordered by descending $J_k$. Define
\begin{equation}\label{eq:ek-def}
e_k = d - J_k - K_k = d - W + J_k.
\end{equation}
Since $W$ is fixed and the $J_k$ are strictly decreasing by the ordering, the $e_k$ are
strictly decreasing, hence distinct. This distinctness is the whole input to the rank lemma
below, so we record it explicitly: for $k < l$,
\[
e_k - e_l = J_k - J_l > 0.
\]

The equations for weight $W$ come from the pairs $(t, i)$ with $2d - t + i = W$, where $t$ is
the $b$-power and $i$ the $\tilde{a}$-power, at both centers. To see why, expand the factor
$(a + \tilde{a} b)^{e_k}$ in \eqref{eq:bridge-pullback} by the binomial theorem, term by
term:
\[
(a + \tilde{a} b)^{e_k} = \sum_{i \geq 0} \binom{e_k}{i} a^{e_k - i} (\tilde{a} b)^i.
\]
Multiplying by $b^{2d - 2J_k - K_k}$ and using $W = 2J_k + K_k$, the monomial $p_{J_k K_k}$
contributes to the coefficient of $b^t$ exactly the terms with $2d - W + i = t$, that is,
$i = t - (2d - W)$. Collecting over $k$ gives the $(t, i, a)$ equation:
\begin{equation}\label{eq:block}
\sum_{k=1}^n p_{J_k K_k}\binom{e_k}{i} a^{e_k-i} = 0, \qquad a \in \{a_3, a_4\}.
\end{equation}
Crucially, the equations \eqref{eq:block} for weight $W$ involve only variables of weight
exactly $W$ at the top: variables of weight $> W$ contribute to higher $b$-powers or have
already been eliminated, as the induction below makes precise.

\begin{lemma}[Block full rank]\label{lem:block-rank}
For distinct $e_1, \dots, e_n$ and $m \geq n$, the $2m \times n$ matrix with entries
$\binom{e_k}{i} a^{e_k-i}$ ($i = 0, \dots, m-1$, $a \in \{a_3, a_4\}$) has rank $n$.
\end{lemma}

\begin{proof}
We reduce the matrix in three explicit steps, each preserving rank.

\emph{Step 1: factor the columns and rows.} Factor column $k$ by $a^{e_k} \neq 0$ (note $a$
is $a_3$ or $a_4$, both nonzero by hypothesis) and row $(i, a)$ by $a^{-i} \neq 0$. Rank is
preserved under nonzero column and row scalings. After this step, the entry in row
$(i, a)$, column $k$ is $\binom{e_k}{i}$, independent of $a$. The two centers now give two
copies of the same rows, so it suffices to show that the $m \times n$ matrix
$[\binom{e_k}{i}]$ ($i = 0, \dots, m-1$) has rank $n$; taking the first $n$ rows
($i = 0, \dots, n-1$) suffices, since $m \geq n$.

\emph{Step 2: triangularize by row operations.} For fixed $i$, the binomial coefficient
$\binom{e}{i}$, viewed as a function of $e$, is a polynomial in $e$ of exact degree $i$:
\[
\binom{e}{i} = \frac{e (e-1) \cdots (e - i + 1)}{i!}
= \frac{1}{i!} e^i + (\text{terms of degree} < i \text{ in } e).
\]
Hence row $i$ of the matrix $[\binom{e_k}{i}]$ equals $\frac{1}{i!}$ times the row
$[e_k^i]$ plus a linear combination of the rows $[e_k^j]$ for $j < i$. Proceeding from
$i = 0$ upward, subtract suitable multiples of earlier rows from row $i$ to kill all
lower-degree terms: row $0$ is already $[1]$; row $1$ is $[e_k] + (\text{constant row})$,
so subtracting the constant multiple of row $0$ leaves $[e_k]$; row $2$ is
$\frac{1}{2}[e_k^2] + (\text{combination of rows } 0, 1)$, so subtracting leaves
$\frac{1}{2}[e_k^2]$; and so on. These are elementary row operations, so they preserve rank
and multiply the determinant by $1$. After all $n$ steps, the matrix is the scaled
Vandermonde matrix $[e_k^i / i!]$ ($i = 0, \dots, n-1$).

\emph{Step 3: the Vandermonde determinant.} The determinant of $[e_k^i]$ is the classical
Vandermonde product $\prod_{j < k} (e_k - e_j)$, nonzero because the $e_k$ are distinct
(shown above from the ordering by $J_k$). Factoring $1/i!$ out of row $i$ gives
\begin{equation}\label{eq:vandermonde-det}
\det[\binom{e_k}{i}] = \frac{\prod_{j<k}(e_k - e_j)}{\prod_{i=0}^{n-1} i!} \neq 0.
\end{equation}
A nonzero $n \times n$ minor means the matrix has rank $n$. The $n = 3$ case is verified
symbolically in \filename{m1\_block\_rank.py}, giving
$-(e_1-e_2)(e_1-e_3)(e_2-e_3)/2$, which agrees with \eqref{eq:vandermonde-det} since
$\prod_{i=0}^{2} i! = 0! \cdot 1! \cdot 2! = 2$ and the sign is the ordering of the factors.
This completes the proof.
\end{proof}

\begin{proof}[Proof of Theorem~\ref{thm:bridge}]
Descend on $W$ from $2d$ to $2$, maintaining the inductive hypothesis that all variables
with weight $> W$ are zero. The base case is $W = 2d$: there are no variables of weight
$> 2d$ (since $W = 2J + K \leq 2d$ for monomials of degree at most $d$), so the hypothesis
holds vacuously.

For the inductive step at weight $W$, the equations \eqref{eq:block} for weight $W$ involve
only variables of weight $\geq W$: a variable of weight $W' > W$ contributes to $b$-powers
above those in the weight-$W$ system, and by the inductive hypothesis all such variables are
already $0$. Hence the weight-$W$ system is exactly the block system of
Lemma~\ref{lem:block-rank}, which has full column rank $n$. Therefore
$p_{J_k K_k} = 0$ for all $k$, that is, all variables of weight $W$ vanish. This restores the
hypothesis for the next lower weight.

Descending to $W = 2$ eliminates every variable with $J \geq 1$: any monomial with $J \geq 1$
has $W = 2J + K \geq 2$, and the descent covers all weights $W \geq 2$. The only monomials
with $W < 2$ are $(J,K) = (0,0)$ with $W = 0$ and $(J,K) = (0,1)$ with $W = 1$. These are never
constrained by the descent, so the nullspace is exactly $\operatorname{span}\{1, y\}$.
\end{proof}

\subsection{Computational verification}

The induction is verified computationally, at the fixed degrees and instances stated (these
are checks of the linear algebra, not replacements for the proof above):
\begin{itemize}[nosep]
\item Symbolic nullspace at general $a_3 \neq a_4$: $d = 2$ gives nullity $2$; $d = 3$ gives
nullity $2$ (\filename{m1\_bridge\_general.py}).
\item Numeric nullspace at $a_3 = 2$, $a_4 = 5$: nullity $2$ for $d = 2, \dots, 8$
(\filename{m1\_induction.py}).
\item Block-triangular structure exhibited at $d = 6$ (\filename{m1\_triangular.py}).
\item Vandermonde determinant verified symbolically (\filename{m1\_block\_rank.py}).
\end{itemize}

\subsection{Scope of the bridge result}\label{sec:bridge-scope}

Under premises (P1) and (P2), the bridge has the following consequences, and no others are
claimed:
\begin{itemize}[nosep]
\item With a single center imposed, the same block argument gives the larger nullspace
$\operatorname{span}\{1, y, x - a y^2\}$ (\filename{lean/TwoDicritical/OneCenter.lean},
\texttt{bridge\_one\_center}). So one such dicritical already forces degree at most $2$:
no pencil of degree at least $3$ has even one dicritical of this chart type.
\item In degree $2$, the automorphism $(x - a y^2, y)$ has exactly one such dicritical,
consistent with Theorem~R (it has one dicritical at infinity, as
Corollary~\ref{cor:automorphisms} requires, and no second).
\item The bridge as proved here says nothing about pencils of degree at least $3$ beyond the
one-center remark above, because (P2) is not derived: a dicritical of degree $1$ with pole
order greater than $1$ is not governed by the system \eqref{eq:block}. Nor does it cover
dicriticals over satellite points, over different points of $L_\infty$, or at different
depths; those are obligation O1.
\end{itemize}

\subsection{The adjacent case: a computation at \texorpdfstring{$d = 3$}{d=3} only}

The adjacent $m = 1$ case --- two dicriticals sharing a curve in the resolution tree rather
than lying over two free points of one $E_2$ --- is treated in this paper only by computation
at $d = 3$ (\filename{e3\_adjacent.py}, \filename{routeB\_cancellation.py}), covering the
$E_2$--$E_3$ adjacent pair at that degree. \textit{Status: grade C, fixed-degree
computation.} There is no argument here for other degrees; the general adjacent case is
obligation O4 in \S\ref{sec:open}. In the October 1 draft, this computation was cited as
closing the adjacent case for every degree. That citation is withdrawn.

\section{What remains open}\label{sec:open}

This section replaces the October 1 draft's ``Proof of the main theorem,'' which assembled
the false $E_1$ theorem, the free-only constant-form sketch, the premises-dependent bridge,
and the $d = 3$ adjacent computation into a contradiction, via a case split (``an adjacent
pair or a bridge'') that was asserted, not proved. No such assembly is attempted here. The
obligations below are what a proof of Theorem~R still needs; the tests after them are checks
any proposed proof must pass.

\subsection{Obligations}

\begin{itemize}[nosep]
\item[\textbf{O1.}] \textbf{The case split.} The positions of two dicriticals in a resolution
tree are more varied than ``an adjacent pair or a bridge'': dicriticals can lie over
\emph{different points} of $L_\infty$ (a pencil can have several base points at infinity in
general, as the non-Keller map $(x, xy)$ shows, and whether a Keller counterexample's pencil
can must be settled); over \emph{satellite points} (the point $a = \infty$ of $E_2$ in
$(x + y^3, y)$, and deeper ones); and at \emph{different depths}. A proof must cover every
configuration of $E_\infty$ and $E_S$ from \S\ref{sec:where}.
\item[\textbf{O2.}] \textbf{Premise (P2).} Derive pole order at most $1$ from the definition of
$m$, or restate the bridge with (P2) as a hypothesis and supply a separate argument for
higher pole order. Until then, the bridge governs only the system
\eqref{eq:bridge-pullback}--\eqref{eq:block} under (P1) and (P2).
\item[\textbf{O3.}] \textbf{$m \geq 2$ at depth $\geq 2$.} Prove the constant-leading-form
statement of Proposition~\ref{prop:constant-form} for every chain of centers, satellite ones
included. Today there is a sketch for free chains only.
\item[\textbf{O4.}] \textbf{The adjacent case} for every degree, not only $d = 3$.
\item[\textbf{O5.}] \textbf{Two small formalizations:} the chart step of Lemma~\ref{lem:e1}
and the binary-form step of Lemma~\ref{lem:linf} in Lean. The algebra of Lemma~\ref{lem:e1}
is formalized; the chart step is executed, not formalized.
\item[\textbf{O6.}] \textbf{The Keller condition must be used globally.} Theorem~R is false
without it. The map $(x, xy)$, with Jacobian $x$, has two dicriticals at infinity in
$\Lambda_{(1,1)}$, over $[1 : 0 : 0]$ and $[0 : 1 : 0]$; the second maps onto the line $u = 0$
of non-proper values (\filename{restatement\_checks.py}, check $2$). Any argument that never
uses $[P,Q] \in K^\times$ cannot prove Theorem~R. This is a standing check on every step
above: Lemmas~\ref{lem:linf} and~\ref{lem:e1} use the Keller condition (vanishing of
$[P_d, Q_d]$, vanishing of the Wronskian coefficient), while the bridge's linear algebra and
the constant-form sketch must be read with their premises attached.
\end{itemize}

\subsection{Tests any proof must pass}

\begin{itemize}[nosep]
\item[\textbf{T1.}] Automorphisms have exactly one dicritical at infinity, of degree $1$
(Corollary~\ref{cor:automorphisms}), but its position varies: over the free point $a$ of
$E_2$ for $(x - a y^2, y)$ (where $P/Q = \tilde{a}/a$ there), and over the satellite point
$a = \infty$ of $E_2$, at depth $5$, for $(x + y^3, y)$. A lemma that rules out a dicritical
in either position is false.
\item[\textbf{T2.}] The identity. Its dicritical at infinity is $L_\infty$ itself;
Lemma~\ref{lem:linf} excludes only non-affine maps.
\item[\textbf{T3.}] Non-Keller maps can have two dicriticals at infinity (O6).
\item[\textbf{T4.}] Every Keller map of degree at least $2$ has two exceptional dicriticals in
total (Corollary~\ref{cor:count-fails}), so only the count \emph{at infinity} can be
bounded. Any argument that bounds the total count is refuted by Example~\ref{ex:automorphism}.
\end{itemize}

\section{Script-to-arithmetic correspondence}\label{sec:correspondence}

Each script below was actually executed; outputs are archived in \filename{outputs/}. The
table maps every script to the exact paper statement it certifies, \emph{under the corrected
statements of this revision}. The grade column follows the restatement note, \S 8: A means
machine-checked, B means a proof by hand from stated facts or an exact computation outside
Lean, C means a sketch or a computation at fixed degrees. A script that certified a withdrawn
statement is listed with what it actually certifies now, not with the draft's reading of it.

\begin{center}
\small
\begin{tabular}{llll}
\textbf{Script} & \textbf{Paper locus} & \textbf{Certifies} & \textbf{Grade} \\ \midrule
\filename{general\_constant\_form.py} & Prop.~\ref{prop:constant-form} & Leading-unit degree count, free chains only & C \\
\filename{middle\_e2.py} & Prop.~\ref{prop:constant-form} & $E_2$ instance of the sketch & C \\
\filename{leaf\_e3.py} & Prop.~\ref{prop:constant-form} & $E_3$ leaf instance of the sketch & C \\
\filename{leaf\_e4.py} & Prop.~\ref{prop:constant-form} & $E_4$ leaf instance of the sketch & C \\
\filename{e1\_solitary.py} & \S\ref{sec:e1} & $12$-monomial count at $d = 4$ (indexing as computed) & C \\
\filename{e1\_global.py} & \S\ref{sec:e1} & Explicit pair, $\nu = -2$, $\deg = 2$ (membership data) & C \\
\filename{e1\_keller\_search.py} & \S\ref{sec:e1} & Searched the misread monomial set; superseded & --- \\
\filename{m1\_bridge\_general.py} & Thm.~\ref{thm:bridge} & Symbolic nullspace, $d = 2, 3$ & B \\
\filename{m1\_induction.py} & \S\ref{sec:bridge} & Nullity $2$ for $d = 2, \dots, 8$ & C \\
\filename{m1\_triangular.py} & \S\ref{sec:bridge} & Block-triangular structure at $d = 6$ & C \\
\filename{m1\_block\_rank.py} & Lem.~\ref{lem:block-rank} & Vandermonde determinant $\neq 0$ & B \\
\filename{at\_degree\_bound.py} & Prop.~\ref{prop:constant-form} & Finite bridge table, $m \leq 4$ & C \\
\filename{at\_degree\_bound\_m5.py} & Prop.~\ref{prop:constant-form} & Finite bridge table, $m \leq 5$ & C \\
\filename{mixed\_poles.py} & Prop.~\ref{prop:constant-form} & Mixed $(2,3)$ bridge instance & C \\
\filename{depth5\_global.py} & Prop.~\ref{prop:constant-form} & Depth-$5$ enumeration & C \\
\filename{depth6\_enum.py} & Prop.~\ref{prop:constant-form} & Depth-$6$ enumeration & C \\
\filename{routeB\_cancellation.py} & \S\ref{sec:bridge} & Adjacent $m = 1$ at $d = 3$ only & C \\
\filename{e3\_adjacent.py} & \S\ref{sec:bridge} & $E_2$--$E_3$ adjacent pair at $d = 3$ & C \\
\end{tabular}
\end{center}

The Lean evidence, which is not a script in the table above, is:
\begin{itemize}[nosep]
\item \filename{lean/TwoDicritical/E1Keller.lean}: \texttt{thm\_e1\_false} (the draft's $E_1$
theorem is false, Example~\ref{ex:e1-counter}); \texttt{e1\_lead\_proportional} and
\texttt{jac2\_eq\_zero\_of\_natDegree\_zero} (the algebra of Lemma~\ref{lem:e1}); grade A for
the algebra.
\item \filename{lean/TwoDicritical/BridgeCollapse.lean} and
\filename{lean/TwoDicritical/OneCenter.lean}: the bridge equations and the one-center span
(Theorem~\ref{thm:bridge}, \S\ref{sec:bridge-scope}); grade A for the equations.
\item \filename{lean/check\_e1\_index.py} (check $5$) and
\filename{lean/check\_satellite.py} (check $7$): the $E_1$ chart step and the
degree-versus-pole-order distinction; executed checks, grade B.
\item \filename{base\_point\_at\_infinity\_check.py} (check $1$) and
\filename{restatement\_checks.py} (check $2$): the binary-form step of
Lemma~\ref{lem:linf} and the non-Keller example of O6; executed checks, grade B.
\end{itemize}

\section{Computational archive}\label{sec:archive}

All scripts used in this work are archived in \filename{scripts/} with their outputs in
\filename{outputs/}, as archived with the October 1 draft:
\begin{itemize}[nosep]
\item 18 Python scripts (SymPy-based), each with a header stating its epistemic status;
\item 18 output files recording exact stdout and exit codes;
\item All scripts exited with code $0$ on re-execution (October 1, 2026).
\end{itemize}
This revision changes no script and claims no new re-execution: the corrections are to the
statements the scripts are read as certifying (\S\ref{sec:correspondence}), not to the
scripts themselves.

\subsection*{Reproducibility}
All artifacts are archived in \href{https://github.com/brandonmccraryresearch-cloud/2D-Jacobian-Conjecture-Workspace}{the public repository}, directory \filename{branch\_two\_dicritical/}. Each script runs with \filename{\textasciitilde/miniconda3/envs/physics/bin/python} (Python 3.12, SymPy). The verification script \filename{verify\_all.sh} re-runs every script and checks exit codes.

\section{Computational and AI-assisted research disclosure}\label{sec:disclosure}

Muse (Meta), the author's personal AI assistant, assisted with computational scripting,
manuscript drafting, and proof verification; all script outputs were verified by actual
execution. The author (B.~D.~McCrary) directed the investigation, verified the results, and
takes responsibility for the manuscript. This revision was drafted with the same assistance:
the October 1 draft was corrected against the October 13, 2026 restatement note, the
derivations were expanded step by step, and no new theorem, computation, or script output
was introduced. The mathematical argument stands on the artifacts above, not on these
attributions.

\begin{thebibliography}{9}
\bibitem{GGHV} Guccione, J.~A., Guccione, J.~J., Horruitiner, R., Valqui, C. \textit{The Jacobian conjecture for the degree pair $(72,108)$.} arXiv:2204.14178v1 (2022).
\bibitem{JC125note} McCrary, B.~D., Muse. \textit{The $(72,108)$ Candidate for the Planar Jacobian Conjecture: A Containment Note.} Zenodo, 2026. \href{https://doi.org/10.5281/zenodo.22941057}{doi:10.5281/zenodo.22941057}.
\bibitem{Jelonek1993} Jelonek, Z. \textit{The set of points at which a polynomial map is not proper.} Ann.\ Polon.\ Math.\ 58 (1993), no.~3, 259--266. doi:10.4064/ap-58-3-259-266.
\bibitem{Chau2008} Nguyen Van Chau. \textit{Plane Jacobian problem for rational polynomials.} arXiv:0801.4138.
\bibitem{Chau2009} Nguyen Van Chau. \textit{Pencil of irreducible rational curves and plane Jacobian conjecture.} arXiv:0905.3939.
\bibitem{AH2015} Abhyankar, S.~S., Heinzer, W.~J. \textit{Existence of dicritical divisors revisited.} arXiv:1508.06015.
\bibitem{AbhyankarSurvey} \textit{Abhyankar's work on dicritical divisors.} arXiv:1707.06733. A survey; it states no equivalence between dicriticals and the conjecture.
\bibitem{BBR1962} Bia\l{}ynicki-Birula, A., Rosenlicht, M. \textit{Injective morphisms of real algebraic varieties.} Proc.\ Amer.\ Math.\ Soc.\ 13 (1962), 200--203.
\bibitem{Restatement} McCrary, B.~D. \textit{The two-dicritical theorem, restated.} Unpublished working note, \filename{branch\_two\_dicritical/RESTATEMENT.md}, October 13, 2026. Not published or deposited.
\end{thebibliography}

\end{document}
